%----------------------------------------------------------------------------------------
% NOTATIE EN FIGUREN VAN DE CURSUS VECTOREN (2de graad)
% Staat in de preambule en niet in een hoofdstuk: ook de lijsten met oplossingen gebruiken dit.
% Vertrekpunt: notatie.tex van Driehoeksmeting (zonder de ruimtefiguren), met het rooster met assen uit Functies.
%----------------------------------------------------------------------------------------
% Bewaarde grafiek, geplaatst door maak-pdf.py: bestand en pagina in figuren-cache/
\newcommand{\grafiekbestand}[2]{\includegraphics[page=#2,trim=3bp 3bp 3bp 3bp]{figuren-cache/#1}}
\usetikzlibrary{angles}

%----------------------------------------------------------------------------------------
% VECTOREN
%----------------------------------------------------------------------------------------
% Een vector, altijd met een pijl erboven: \vect{v} voor een vector met een naam, \vect{AB} voor de
% vector van A naar B. Dezelfde pijl in beide gevallen, zodat de notatie overal gelijk is.
\newcommand{\vect}[1]{\overrightarrow{#1}}
% De nulvector
\newcommand{\nulvec}{\vect{0}}
% De norm (de lengte) van een vector: \norm{\vect{v}} geeft ||v||
\newcommand{\norm}[1]{\lVert #1\rVert}
% Componenten van een vector: \vc{4}{-3} geeft (4, -3). De accolades houden een minteken bij het getal,
% en de komma is een scheidingsteken, geen decimale komma.
\newcommand{\vc}[2]{({#1}, {#2})}

% Lengte van een lijnstuk: \lengte{AB} geeft |AB|; \abs{-2} de absolute waarde.
% Met \lvert en \rvert, zodat een minteken na de streep geen aftrekking wordt.
\newcommand{\lengte}[1]{\lvert #1\rvert}
\newcommand{\abs}[1]{\lvert #1\rvert}
% Driehoek en hoek: \drieh{ABC} geeft de driehoek ABC, \hoek{A} een hoek met een dakje
\newcommand{\drieh}[1]{\triangle #1}
\newcommand{\hoek}[1]{\widehat{#1}}

% Zoals stappen, maar met een bredere kolom voor de berekening: voor lange formules
\newenvironment{stappenbreed}{%
  \par\renewcommand{\arraystretch}{1.25}%
  \noindent\begin{tabular}{@{}>{$\displaystyle}p{9.8cm}<{$}@{\hspace{4mm}}>{\RaggedRight\small\itshape\arraybackslash}p{4.9cm}@{}}%
}{\end{tabular}\par}

% Zoals stappen, met eigen kolombreedtes: \begin{stappenvrij}{6.2cm}{3.8cm}. Voor naast een figuur.
\newenvironment{stappenvrij}[2]{%
  \par\renewcommand{\arraystretch}{1.25}%
  \noindent\begin{tabular}{@{}>{$\displaystyle}p{#1}<{$}@{\hspace{3mm}}>{\RaggedRight\small\itshape\arraybackslash}p{#2}@{}}%
}{\end{tabular}\par}

% Tekst of berekening links, figuur rechts: \naast{links}{rechts}
\newsavebox{\naastbox}
\newcommand{\naast}[3][]{%
  \par\noindent\sbox{\naastbox}{#3}%
  \begin{minipage}[t]{\dimexpr\linewidth-\wd\naastbox-3mm\relax}\vspace{0pt}#2\end{minipage}\hfill
  \begin{minipage}[t]{\wd\naastbox}\vspace{0pt}\usebox{\naastbox}\end{minipage}\par\addvspace{1mm}}
% Een rij figuren, gecentreerd
\newcommand{\figuren}[1]{\par\addvspace{1mm}{\centering #1\par}\addvspace{1.5mm}}
% Opgave met een figuur rechts: \begin{metgrafiek}[breedte]{tikz} opgave en deelvragen \end{metgrafiek}
\newsavebox{\grafiekbox}
\newenvironment{metgrafiek}[2][]{%
  \sbox{\grafiekbox}{#2}%
  \begin{minipage}[t]{\dimexpr\linewidth-\wd\grafiekbox-3mm\relax}%
}{%
  \end{minipage}\hfill
  \begin{minipage}[t]{\wd\grafiekbox}\vspace{-2mm}\usebox{\grafiekbox}\end{minipage}\par}

%----------------------------------------------------------------------------------------
% FIGUREN, binnen een tikzpicture
%----------------------------------------------------------------------------------------
% Pijlen voor vectoren. vector: een gegeven vector (kleur van het hoofdstuk); vectorb: een tweede
% gegeven vector (donker); som: de uitkomst (dik, kleur van het hoofdstuk); hulpvector: een verschoven
% kopie of een tussenstap (grijs, gestreept). Zo verschillen ze ook in zwart-wit.
\tikzset{
  vector/.style={-{Stealth[length=2.6mm,width=2mm]},line width=1.1pt,accent},
  vectorb/.style={vector,inkt},
  som/.style={-{Stealth[length=3mm,width=2.4mm]},line width=1.8pt,accent},
  hulpvector/.style={-{Stealth[length=2.4mm,width=1.8mm]},line width=.9pt,grijs,densely dashed},
  hulplijn/.style={grijs,thin,densely dashed}}
% Een pijl: \pijl[stijl]{(x1,y1)}{(x2,y2)}; standaard de stijl vector
\newcommand{\pijl}[3][vector]{\draw[#1] #2 -- #3;}

% Rooster zonder assen: \ruitjes{xmin}{ymin}{xmax}{ymax}
\newcommand{\ruitjes}[4]{\draw[black!18,line width=.3pt] (#1,#2) grid (#3,#4);}
% Rooster met assen (zoals in Functies): \rooster{xmin}{xstap}{xmax}{ymin}{ystap}{ymax}, met getallen
% bij elke lijn (niet bij 0). \roosterleeg doet hetzelfde zonder getallen; getallen bij een grotere stap
% met \xgetallen{2,4} en \ygetallen{2,4}.
\NewDocumentCommand{\roosterleeg}{O{$x$} O{$y$} m m m m m m}{%
  \draw[black!14,very thin,xstep=#4,ystep=#7] (#3,#6) grid (#5,#8);
  \draw[->,grijs] (#3,0) -- ({#5+.45*#4},0) node[right,text=black,font=\footnotesize] {#1};
  \draw[->,grijs] (0,#6) -- (0,{#8+.45*#7}) node[above,text=black,font=\footnotesize] {#2};}
\NewDocumentCommand{\rooster}{O{$x$} O{$y$} m m m m m m}{%
  \roosterleeg[#1][#2]{#3}{#4}{#5}{#6}{#7}{#8}
  \pgfmathsetmacro{\rtweex}{#3+#4}\pgfmathsetmacro{\rtweey}{#6+#7}
  \foreach \t in {#3,\rtweex,...,#5}{\pgfmathtruncatemacro{\tt}{round(\t)}\ifdim\t pt=0pt\else\node[below,inner sep=1.5pt,text=black,font=\scriptsize] at (\t,0) {$\tt$};\fi}
  \foreach \t in {#6,\rtweey,...,#8}{\pgfmathtruncatemacro{\tt}{round(\t)}\ifdim\t pt=0pt\else\node[left,inner sep=1.5pt,text=black,font=\scriptsize] at (0,\t) {$\tt$};\fi}}
\newcommand{\xgetallen}[1]{\foreach \t in {#1}{\node[below,inner sep=1.5pt,text=black,font=\scriptsize] at (\t,0) {$\t$};}}
\newcommand{\ygetallen}[1]{\foreach \t in {#1}{\node[left,inner sep=1.5pt,text=black,font=\scriptsize] at (0,\t) {$\t$};}}
% De oorsprong: \oorsprong zet een 0 linksonder de oorsprong
\newcommand{\oorsprong}{\node[below left,inner sep=1.5pt,text=black,font=\scriptsize] at (0,0) {$0$};}

% Een punt als stip: \stip{x}{y}, of met label: \punt[positie]{x}{y}{label}
\newcommand{\stip}[2]{\fill (#1,#2) circle (1.5pt);}
\newcommand{\punt}[4][above]{\fill (#2,#3) circle (1.5pt) node[#1,inner sep=2pt] {$#4$};}
% Een hoekboogje bij hoekpunt B tussen de stralen naar A en C (tegenwijzerzin van A naar C):
% \hoekboog[opties]{A}{B}{C}{straal}; werkt met benoemde coördinaten
\newcommand{\hoekboog}[5][]{\pic[draw,accent,angle radius=#5,#1] {angle=#2--#3--#4};}
% Hoeklabel (zonder de bibliotheek quotes, die botst met de aanhalingstekens van babel): \hoeklabel[opties]{A}{B}{C}{straal}{tekst}: boogje bij B en de tekst in de hoek
\newcommand{\hoeklabel}[6][]{\pic[draw,accent,angle radius=#5,pic text={$#6$},angle eccentricity=1.55,pic text options={font=\scriptsize},#1] {angle=#2--#3--#4};}
% Een rechte hoek bij B tussen A en C: \rechtehoek{A}{B}{C}
\newcommand{\rechtehoek}[3]{\pic[draw,grijs,angle radius=2.2mm] {right angle=#1--#2--#3};}
% Hetzelfde met coördinaten: \rechtehoekxy{(0,3)}{(0,0)}{(4,0)}
\newcommand{\rechtehoekxy}[3]{\coordinate (rh@a) at #1; \coordinate (rh@b) at #2; \coordinate (rh@c) at #3;
  \pic[draw,grijs,angle radius=2.2mm] {right angle=rh@a--rh@b--rh@c};}
